\documentclass{article}
\usepackage{amsmath,amssymb,amsthm}
\usepackage{algorithm}
\usepackage{algpseudocode}
\usepackage{enumitem}
\usepackage{hyperref}
\usepackage{geometry}
\geometry{margin=1in}
\newtheorem{theorem}{Theorem}
\newtheorem{lemma}{Lemma}
\newtheorem{assumption}{Assumption}
\newcommand{\prox}{\operatorname{prox}}
\newcommand{\grad}{\nabla}
\newcommand{\R}{\mathbb{R}}

\begin{document}

\section*{Algorithm Name}
\textbf{Proximal Gradient Method (PGM) for Non‑Convex Composite Optimization}  

We consider the composite problem  
\[
\min_{x\in\R^{d}} \; F(x)\;:=\; f(x)+g(x),
\]
where $f$ is smooth (possibly non‑convex) and $g$ is proper, lower‑semicontinuous, possibly non‑smooth but with an efficiently computable proximal operator.

\section*{Mathematical Formulation}
Given a stepsize $\eta>0$, the iterates $\{x_k\}_{k\ge 0}$ are generated by  
\begin{equation}
\boxed{%
x_{k+1}= \prox_{\eta g}\!\bigl(x_k-\eta\,\grad f(x_k)\bigr)},
\qquad k=0,1,2,\dots
\tag{1}
\end{equation}
where for any $z\in\R^{d}$  
\[
\prox_{\eta g}(z):=\arg\min_{u\in\R^{d}}\Bigl\{ g(u)+\frac{1}{2\eta}\|u-z\|^{2}\Bigr\}.
\]

Define the \emph{proximal gradient mapping}
\[
G_{\eta}(x):=\frac{1}{\eta}\Bigl(x-\prox_{\eta g}(x-\eta\grad f(x))\Bigr).
\tag{2}
\]
When $g\equiv0$, $G_{\eta}(x)=\grad f(x)$. A point $x^{\star}$ satisfying $0\in\grad f(x^{\star})+\partial g(x^{\star})$ is a stationary point of $F$; equivalently $G_{\eta}(x^{\star})=0$.

\section*{Pseudocode}
\begin{algorithm}[H]
\caption{Proximal Gradient Method}
\begin{algorithmic}[1]
\State \textbf{Input:} stepsize $\eta>0$, max iterations $K$, tolerance $\varepsilon>0$.
\State Initialize $x_{0}\in\R^{d}$.
\For{$k=0,1,\dots,K-1$}
    \State $v_k \gets \grad f(x_k)$                     \Comment{gradient of smooth part}
    \State $y_k \gets x_k - \eta\, v_k$                \Comment{gradient step}
    \State $x_{k+1} \gets \prox_{\eta g}(y_k)$         \Comment{proximal step}
    \State Compute $G_{\eta}(x_k)=\frac{1}{\eta}(x_k-x_{k+1})$.
    \If{$\|G_{\eta}(x_k)\|\le \varepsilon$}
        \State \textbf{break}
    \EndIf
\EndFor
\State \textbf{Output:} $x_{k+1}$.
\end{algorithmic}
\end{algorithm}

\section*{Dry Run Example}
Consider the simple scalar problem  
\[
f(w)=(w-3)^{2},\qquad g\equiv0,\qquad \eta=0.1,\qquad w_{0}=0.
\]
Since $g\equiv0$, $\prox_{\eta g}= \operatorname{id}$.

\begin{center}
\begin{tabular}{c|c|c|c}
Iteration $k$ & $w_k$ & $\grad f(w_k)=2(w_k-3)$ & $w_{k+1}=w_k-\eta\,\grad f(w_k)$ \\ \hline
0 & $0$ & $-6$ & $0-0.1(-6)=0.6$ \\
1 & $0.6$ & $2(0.6-3)=-4.8$ & $0.6-0.1(-4.8)=1.08$ \\
2 & $1.08$ & $2(1.08-3)=-3.84$ & $1.08-0.1(-3.84)=1.464$ \\
3 & $1.464$ & $2(1.464-3)=-3.072$ & $1.464-0.1(-3.072)=1.7712$ 
\end{tabular}
\end{center}
The objective values $f(w_k)$ decrease from $9$ to $2.95$, illustrating the descent property.

\newpage
\section*{Assumptions}
\begin{assumption}[Smoothness of $f$]\label{ass:smooth}
$f:\R^{d}\to\R$ is differentiable and its gradient is $L$–Lipschitz:
\[
\|\grad f(x)-\grad f(y)\|\;\le\;L\|x-y\|,\qquad\forall x,y\in\R^{d}.
\tag{A1}
\]
\end{assumption}
\begin{assumption}[Bounded Below]\label{ass:bound}
The composite objective $F(x)=f(x)+g(x)$ is bounded from below:
\[
F^{\inf}:=\inf_{x\in\R^{d}}F(x)>-\infty .
\tag{A2}
\end{assumption}
\begin{assumption}[Stepsize]\label{ass:stepsize}
The stepsize satisfies
\[
0<\eta\le\frac{1}{L}.
\tag{A3}
\end{assumption}
\begin{assumption}[Proximal Operator]\label{ass:prox}
For any $\eta>0$, $\prox_{\eta g}$ is \emph{non‑expansive}:
\[
\|\prox_{\eta g}(u)-\prox_{\eta g}(v)\|\;\le\;\|u-v\|,\qquad\forall u,v\in\R^{d}.
\tag{A4}
\]
\end{assumption}

\section*{Main Convergence Theorem}
\begin{theorem}[Convergence of PGM under non‑convex $f$]\label{thm:main}
Let Assumptions~\ref{ass:smooth}–\ref{ass:prox} hold and generate $\{x_k\}$ by \eqref{1} with a stepsize $\eta$ satisfying Assumption~\ref{ass:stepsize}. Then for any integer $K\ge1$,
\[
\boxed{%
\min_{0\le k\le K-1}\|G_{\eta}(x_k)\|^{2}
\;\le\;
\frac{2\bigl(F(x_{0})-F^{\inf}\bigr)}{\eta\,K}}.
\tag{3}
\]
Consequently,
\[
\lim_{K\to\infty}\min_{0\le k\le K-1}\|G_{\eta}(x_k)\|=0,
\]
i.e. every limit point of $\{x_k\}$ is a stationary point of $F$.
\end{theorem}

\section*{Theoretical Properties}
\begin{itemize}[leftmargin=*]
\item \textbf{Descent guarantee.} Equation \eqref{3} is derived from a per‑iteration decrease of $F$ proportional to $\|G_{\eta}(x_k)\|^{2}$.
\item \textbf{Stepsize sensitivity.} The bound requires $\eta\le1/L$; larger $\eta$ may violate the descent Lemma~\ref{lem:descent}.
\item \textbf{Comparison.} When $g\equiv0$, $G_{\eta}(x)=\grad f(x)$ and \eqref{3} reduces to the classical $O(1/K)$ guarantee for gradient descent on smooth non‑convex functions.
\item \textbf{Stability.} Because $\prox_{\eta g}$ is non‑expansive (Assumption~\ref{ass:prox}), the iterates cannot diverge as long as $F$ is bounded below.
\end{itemize}

\section*{Discussion}
The theorem shows that the proximal gradient mapping, which measures first‑order stationarity of the composite objective, vanishes at a rate $O(1/K)$. No convexity of $f$ is required; only smoothness and a stepsize not exceeding $1/L$. The proof hinges on two elementary lemmas: the smoothness (descent) lemma for $f$ and the non‑expansiveness of the proximal operator. Both are applied in a straightforward algebraic manipulation that yields a telescoping sum.

\newpage
\section*{Preliminaries (Definitions and Lemmas)}
\begin{lemma}[Descent Lemma for $L$‑smooth $f$]\label{lem:descent}
For any $x,y\in\R^{d}$,
\[
f(y)\le f(x)+\langle\grad f(x),y-x\rangle+\frac{L}{2}\|y-x\|^{2}.
\tag{4}
\]
\end{lemma}
\begin{proof}
Standard consequence of $L$‑Lipschitz gradient; see e.g. Nesterov (2004). \qedhere
\end{proof}

\begin{lemma}[Non‑expansiveness of the proximal operator]\label{lem:nonexp}
For any $u,v\in\R^{d}$ and any $\eta>0$,
\[
\bigl\|\prox_{\eta g}(u)-\prox_{\eta g}(v)\bigr\|
\le \|u-v\|.
\tag{5}
\]
\end{lemma}
\begin{proof}
Follows from the firmly non‑expansive property of the proximal mapping; see e.g. Parikh \& Boyd (2014). \qedhere
\end{proof}

\section*{Main Proof (Step‑by‑Step Derivation)}
We prove Theorem~\ref{thm:main}. All steps are numbered for easy reference.

\medskip
\noindent\textbf{Step 1 (Notation).} Define $y_k:=x_k-\eta\grad f(x_k)$. By definition of the algorithm,
\[
x_{k+1}= \prox_{\eta g}(y_k). \tag{6}
\]

\noindent\textbf{Step 2 (Optimality condition of the proximal step).} The definition of $\prox$ yields
\[
0\in \partial g(x_{k+1})+\frac{1}{\eta}(x_{k+1}-y_k),
\]
or equivalently
\[
\frac{1}{\eta}(y_k-x_{k+1})\in \partial g(x_{k+1}). \tag{7}
\]

\noindent\textbf{Step 3 (Expression of the proximal gradient mapping).} Using \eqref{6} and \eqref{7},
\[
G_{\eta}(x_k)=\frac{1}{\eta}(x_k-x_{k+1})
               =\grad f(x_k)+\underbrace{\frac{1}{\eta}(y_k-x_{k+1})}_{\in\partial g(x_{k+1})}.
\tag{8}
\]
Thus $-G_{\eta}(x_k)\in \grad f(x_k)+\partial g(x_{k+1})$.

\noindent\textbf{Step 4 (Descent of $f$).} Apply Lemma~\ref{lem:descent} with $x=x_k$, $y=x_{k+1}$:
\[
f(x_{k+1})\le f(x_k)+\langle\grad f(x_k),x_{k+1}-x_k\rangle
               +\frac{L}{2}\|x_{k+1}-x_k\|^{2}. \tag{9}
\]

\noindent\textbf{Step 5 (Combine with the proximal step).} Using \eqref{8},
\[
\langle\grad f(x_k),x_{k+1}-x_k\rangle
= -\langle G_{\eta}(x_k),x_{k+1}-x_k\rangle
  +\Big\langle\frac{1}{\eta}(y_k-x_{k+1}),x_{k+1}-x_k\Big\rangle .
\tag{10}
\]
Observe that $y_k-x_{k+1}=x_k-\eta\grad f(x_k)-x_{k+1}=-(x_{k+1}-x_k)-\eta\grad f(x_k)$. Substituting,
\[
\Big\langle\frac{1}{\eta}(y_k-x_{k+1}),x_{k+1}-x_k\Big\rangle
= -\frac{1}{\eta}\|x_{k+1}-x_k\|^{2}
  -\langle\grad f(x_k),x_{k+1}-x_k\rangle .
\tag{11}
\]
Insert \eqref{11} into \eqref{10}:
\[
\langle\grad f(x_k),x_{k+1}-x_k\rangle
= -\langle G_{\eta}(x_k),x_{k+1}-x_k\rangle
   -\frac{1}{\eta}\|x_{k+1}-x_k\|^{2}
   -\langle\grad f(x_k),x_{k+1}-x_k\rangle .
\]
Bring the term $ \langle\grad f(x_k),x_{k+1}-x_k\rangle$ to the left side, we obtain
\[
2\langle\grad f(x_k),x_{k+1}-x_k\rangle
   = -\langle G_{\eta}(x_k),x_{k+1}-x_k\rangle
     -\frac{1}{\eta}\|x_{k+1}-x_k\|^{2}.
\tag{12}
\]
Hence
\[
\langle\grad f(x_k),x_{k+1}-x_k\rangle
   = -\frac{1}{2}\langle G_{\eta}(x_k),x_{k+1}-x_k\rangle
     -\frac{1}{2\eta}\|x_{k+1}-x_k\|^{2}.
\tag{13}
\]

\noindent\textbf{Step 6 (Plug back into the descent of $f$).} Replace the inner product in \eqref{9} using \eqref{13}:
\[
\begin{aligned}
f(x_{k+1}) &\le f(x_k)
            -\frac{1}{2}\langle G_{\eta}(x_k),x_{k+1}-x_k\rangle
            -\frac{1}{2\eta}\|x_{k+1}-x_k\|^{2}
            +\frac{L}{2}\|x_{k+1}-x_k\|^{2}\\[4pt]
          &= f(x_k)
            -\frac{1}{2}\langle G_{\eta}(x_k),x_{k+1}-x_k\rangle
            +\frac{1}{2}\Bigl(L-\frac{1}{\eta}\Bigr)\|x_{k+1}-x_k\|^{2}.
\end{aligned}
\tag{14}
\]

\noindent\textbf{Step 7 (Bound the inner product).} By Cauchy–Schwarz and the definition of $G_{\eta}$,
\[
\langle G_{\eta}(x_k),x_{k+1}-x_k\rangle
= -\eta\|G_{\eta}(x_k)\|^{2},
\tag{15}
\]
because $x_{k+1}=x_k-\eta G_{\eta}(x_k)$ from \eqref{2}. Substituting \eqref{15} into \eqref{14} gives
\[
f(x_{k+1})\le f(x_k)
           -\frac{\eta}{2}\|G_{\eta}(x_k)\|^{2}
           +\frac{1}{2}\Bigl(L-\frac{1}{\eta}\Bigr)\|x_{k+1}-x_k\|^{2}.
\tag{16}
\]

\noindent\textbf{Step 8 (Use the stepsize condition).} By Assumption~\ref{ass:stepsize}, $\eta\le 1/L$, i.e. $L-\frac{1}{\eta}\le 0$. Hence the last term in \eqref{16} is non‑positive and can be dropped:
\[
f(x_{k+1})\le f(x_k)-\frac{\eta}{2}\|G_{\eta}(x_k)\|^{2}.
\tag{17}
\]

\noindent\textbf{Step 9 (Add the nonsmooth part).} Since $x_{k+1}$ minimizes the proximal subproblem,
\[
g(x_{k+1})+\frac{1}{2\eta}\|x_{k+1}-y_k\|^{2}
\le g(x_k)+\frac{1}{2\eta}\|x_k-y_k\|^{2}.
\tag{18}
\]
Recall $y_k=x_k-\eta\grad f(x_k)$, thus $\|x_k-y_k\|^{2}= \eta^{2}\|\grad f(x_k)\|^{2}$ and
$\|x_{k+1}-y_k\|^{2}= \eta^{2}\|G_{\eta}(x_k)\|^{2}$ (again using \eqref{2}). Multiplying \eqref{18} by $\frac{1}{2\eta}$ and rearranging yields
\[
g(x_{k+1})\le g(x_k)-\frac{\eta}{2}\|G_{\eta}(x_k)\|^{2}.
\tag{19}
\]

\noindent\textbf{Step 10 (Combine smooth and nonsmooth descent).} Adding \eqref{17} and \eqref{19}:
\[
F(x_{k+1}) = f(x_{k+1})+g(x_{k+1})
\le F(x_k)-\eta\|G_{\eta}(x_k)\|^{2}.
\tag{20}
\]

\noindent\textbf{Step 11 (Telescoping sum).} Summing \eqref{20} for $k=0,\dots,K-1$ we obtain
\[
F(x_{K}) \le F(x_{0})-\eta\sum_{k=0}^{K-1}\|G_{\eta}(x_k)\|^{2}.
\tag{21}
\]
Since $F(x_{K})\ge F^{\inf}$ by Assumption~\ref{ass:bound},
\[
\eta\sum_{k=0}^{K-1}\|G_{\eta}(x_k)\|^{2}
\le F(x_{0})-F^{\inf}.
\tag{22}
\]

\noindent\textbf{Step 12 (Deriving the rate).} Divide \eqref{22} by $K$ and use the definition of the minimum over the first $K$ iterates:
\[
\min_{0\le k\le K-1}\|G_{\eta}(x_k)\|^{2}
\le \frac{1}{K}\sum_{k=0}^{K-1}\|G_{\eta}(x_k)\|^{2}
\le \frac{F(x_{0})-F^{\inf}}{\eta\,K}.
\tag{23}
\]
Multiplying numerator and denominator by $2$ yields exactly the bound \eqref{3} stated in Theorem~\ref{thm:main}. \qed

\section*{Rate Derivation (With Constants)}
From \eqref{23} we have the explicit constant
\[
C:=\frac{2\bigl(F(x_{0})-F^{\inf}\bigr)}{\eta},
\]
so that for any $K\ge1$,
\[
\min_{0\le k\le K-1}\|G_{\eta}(x_k)\|^{2}\le \frac{C}{K}=O\!\left(\frac{1}{K}\right).
\]
If one wishes to bound the \emph{average} squared norm instead of the minimum, (22) directly gives
\[
\frac{1}{K}\sum_{k=0}^{K-1}\|G_{\eta}(x_k)\|^{2}\le \frac{F(x_{0})-F^{\inf}}{\eta\,K}=O\!\left(\frac{1}{K}\right).
\]

\section*{Special Cases}
\begin{enumerate}
\item \textbf{$g\equiv0$ (pure gradient descent).} Then $G_{\eta}(x)=\grad f(x)$ and Theorem~\ref{thm:main} reduces to the classical $O(1/K)$ guarantee for smooth non‑convex optimization.
\item \textbf{Indicator of a convex set $\mathcal{C}$.} If $g=\iota_{\mathcal{C}}$, then $\prox_{\eta g}$ is the Euclidean projection onto $\mathcal{C}$. The algorithm becomes projected gradient descent, and the same $O(1/K)$ bound holds for the projected gradient mapping.
\item \textbf{Strongly convex $g$ with parameter $\mu>0$.} Although the theorem does not require convexity, if $g$ is $\mu$‑strongly convex one can strengthen \eqref{20} to obtain linear convergence under additional error‑bound conditions; this lies beyond the present scope.
\end{enumerate}

\end{document}